\documentclass[11pt,a4paper]{article}
\usepackage[margin=0.85in]{geometry}
\usepackage{graphicx}
\usepackage{booktabs}
\usepackage{hyperref}
\hypersetup{
  colorlinks=true,
  linkcolor=black,
  citecolor=black,
  urlcolor=blue,
  pdfborder={0 0 0}
}
\usepackage{amsmath}
\usepackage{caption}
\usepackage{float}
\usepackage{placeins}
\usepackage{setspace}
\usepackage{microtype}
\usepackage{fancyhdr}
\usepackage{titlesec}
\usepackage{enumitem}
\setstretch{1.12}
\graphicspath{{outputs/}}
\setlist[itemize]{leftmargin=1.4em,itemsep=0.25em,topsep=0.35em}

\captionsetup{font=small,labelfont=bf,skip=4pt}
\setlength{\textfloatsep}{8pt plus 2pt minus 2pt}
\setlength{\floatsep}{6pt plus 2pt minus 2pt}
\setlength{\intextsep}{8pt plus 2pt minus 2pt}
\titleformat{\section}{\large\bfseries}{\thesection}{0.6em}{}
\titlespacing*{\section}{0pt}{1.1ex plus 0.3ex}{0.6ex}
\titleformat{\paragraph}[runin]{\bfseries}{}{0pt}{}[.~]
\titlespacing*{\paragraph}{0pt}{0.8ex}{0.5em}

\newcommand{\safeimg}[2][0.88\linewidth]{%
  \includegraphics[width=#1,height=0.40\textheight,keepaspectratio]{#2}%
}

\pagestyle{fancy}
\fancyhf{}
\fancyhead[L]{\footnotesize CISC3024 AI Assignment \#1}
\fancyhead[R]{\footnotesize CHE CHI HIN, Angus (UC325182)}
\fancyfoot[C]{\thepage}
\renewcommand{\headrulewidth}{0.3pt}

\title{\textbf{AI Assignment \#1}\\[0.25em]
\large FasterNet for Robust Digit Pattern Recognition\\
under Geometric and Noise Corruptions}
\author{CHE CHI HIN, Angus\\Student ID: UC325182\\CISC3024 Pattern Recognition}
\date{}

\begin{document}
\maketitle
\thispagestyle{fancy}

\begin{abstract}
\textbf{Thesis:} a recent efficient Deep CNN (FasterNet with Partial Convolution) plus robust augmentation and rotation TTA recognizes digits under geometric/noise corruptions far better than clean-only training, and remains competitive with a student-scale equivariant baseline while staying easy for an AI agent to implement and verify.
This report follows the required structure: (1)~how AI was asked to find the algorithm, (2)~algorithm description, (3)~how AI implemented it, (4)~experiment settings and results, (5)~what was learnt, and (6)~the source-code link. We evaluate MNIST, EMNIST Digits and SVHN. Search, coding, experiments and drafting were done by a Cursor AI agent under my direction (no hand-written code).
\end{abstract}

\section{How I Asked AI Tools to Find the Algorithm}
I used \textbf{Cursor} as the only AI tool for search, coding and writing; I did not hand-write code. Concretely, I first asked the agent to \emph{plan and ask clarifying questions} rather than jump into a week-long schedule, because the Moodle brief was short. I then locked constraints in plain language:
\begin{itemize}
  \item Find a \textbf{recent} (2023--2024) Deep CNN/Autoencoder with a peer-reviewed paper; do not copy a peer's \textbf{RepViT} + $8{\times}8$ fusion story.
  \item Prefer an operator we can \textbf{verify} (not accuracy-only), and frame the application as \textbf{robust digit recognition} (a rotated ``3'' should still be recognized).
  \item After early demos looked ``only digit 3 / only $90^\circ$'', I told the agent to enrich multi-digit, multi-angle evidence and to extend to EMNIST/SVHN.
\end{itemize}
The agent proposed several lightweight CNNs; I chose \textbf{FasterNet} (Chen et al., CVPR~2023) for Partial Convolution (PConv), which is distinct from re-parameterization. Diary: \texttt{process\_log.md}.

\section{Algorithm Description}
\paragraph{FasterNet / PConv.}
FasterNet argues that reducing FLOPs alone (e.g., depthwise convolution) may not maximize latency because of memory access cost~\cite{chen2023fasternet}. \textbf{Partial Convolution} applies a spatial $3{\times}3$ convolution to only a fraction $r{=}1/n_{\mathrm{div}}$ of channels (paper default $r{=}1/4$) and leaves the remaining channels untouched; a residual pointwise MLP then mixes channels.

\paragraph{Our adaptation (FasterNet-mini).}
For $28{\times}28$ greyscale (MNIST/EMNIST) and $32{\times}32$ RGB (SVHN) we use a $2{\times}2$ patch stem and three stages (embed 32, depths $(1,2,2)$). Each block is $\mathrm{PConv}(n_{\mathrm{div}})\!\rightarrow\!$ BN+activation pointwise MLP. Main runs use ReLU; we ablate GELU. We also implement a student-scale \textbf{p4 orbit-convolution} baseline for literature comparison (augmentation vs architectural invariance)~\cite{rath2023invariances}.

\begin{figure}[H]
\centering
\safeimg[0.88\linewidth]{fig_architecture.png}
\caption{FasterNet-mini architecture used for MNIST/EMNIST.}
\label{fig:arch}
\end{figure}

\paragraph{PConv verification.}
Untouched channels remain exact identities (max abs error $=0$). At $n_{\mathrm{div}}{=}4$, PConv uses $6.25\%$ of dense $3{\times}3$ FLOPs on a fixed activation shape.

\paragraph{Related context.}
MobileOne~\cite{vasu2023mobileone} targets mobile latency via re-parameterization (different mechanism). Corruption-style evaluation follows Hendrycks \& Dietterich~\cite{hendrycks2019benchmarking}. Digit domains: MNIST~\cite{lecun1998gradient}, EMNIST Digits~\cite{cohen2017emnist}, SVHN~\cite{netzer2011reading}.

\section{How AI Implemented the Algorithm}
Under my direction, the Cursor agent implemented the full pipeline with \textbf{no hand-written student code}:
\begin{itemize}
  \item \texttt{models/fasternet\_mini.py} --- PConv, MLPBlock, FasterNet-mini / T0-lite, GELU option.
  \item \texttt{models/equivariant\_cnn.py} --- p4 orbit-CNN baseline; \texttt{models/baseline\_cnn.py} --- SimpleCNN control.
  \item \texttt{corruptions.py} --- train-time robust aug + test corruptions (rotate/noise/blur/shift/occlusion/elastic).
  \item \texttt{train.py} / \texttt{evaluate.py} --- MNIST/EMNIST/SVHN loaders, MPS training, Rot90 TTA, reject curves, figures JSON.
  \item Analysis: \texttt{scripts/pconv\_efficiency.py}, \texttt{scripts/analyze\_extra.py}, \texttt{scripts/make\_rich\_figures.py}.
  \item Drivers: \texttt{run\_all.py}, \texttt{scripts/run\_upgrade\_15.py}, \texttt{scripts/run\_upgrade\_gelu\_eq\_svhn.py}.
  \item Report/demo: \texttt{report.tex}, \texttt{scripts/app\_gradio.py}, \texttt{process\_log.md}.
\end{itemize}
I set goals, approved FasterNet, requested EMNIST/SVHN, multi-seed stats, GELU/equivariant/longer SVHN, and PDF polish; the agent wrote, ran and iterated the code.

\section{Experiment Settings and Results}
\textbf{Settings.} AdamW ($10^{-3}$), cosine schedule, batch 128, Apple MPS; main MNIST/EMNIST runs use 12 epochs (SVHN ablation: 12 vs 24). Robust-aug includes cardinal/continuous rotations and noise/blur/shift/occlusion/elastic. EMNIST/SVHN use 60k/10k student-scale subsamples. Metrics: corruption accuracy, Rot90+TTA, reject@0.7, ECE, multi-seed mean$\pm$std; ablations on $n_{\mathrm{div}}$, GELU and equivariant baseline.

\begin{table}[H]
\centering
\small
\caption{Main robustness results.}
\label{tab:main}
\begin{tabular}{lccccc}
\toprule
Setting & Clean & Rot45 & Rot90 & Rot90+TTA & Noise \\
\midrule
MNIST FasterNet-aug & 0.987 & 0.966 & 0.887 & 0.948 & 0.891 \\
EMNIST FasterNet-aug & 0.988 & 0.968 & 0.835 & 0.941 & 0.801 \\
MNIST clean-only & 0.993 & 0.588 & 0.127 & 0.659 & 0.377 \\
MNIST Baseline-aug & 0.958 & 0.735 & 0.267 & -- & 0.934 \\
SVHN 12ep & 0.789 & 0.576 & 0.572 & 0.720 & 0.641 \\
SVHN 24ep long & 0.858 & 0.711 & 0.671 & 0.792 & 0.693 \\
GELU (MNIST) & 0.988 & 0.973 & 0.866 & 0.935 & 0.857 \\
p4 Equivariant-aug & 0.964 & 0.912 & 0.823 & 0.927 & 0.831 \\
\bottomrule
\end{tabular}
\end{table}

{\small
Multi-seed MNIST Rot90: $84.1\%\pm 10.0\%$. EMNIST multi-seed Rot90: $81.2\%\pm 3.2\%$. Clean ECE $\approx 0.0071$; hardest Rot90 pair $6\leftrightarrow 9$. Full $n_{\mathrm{div}}\!\in\!\{2,4,8\}$ 12-epoch ablation (Table~\ref{tab:ndiv}) supports paper default $r{=}1/4$ for latency/clean trade-off.
}

\begin{table}[H]
\centering
\small
\caption{$n_{\mathrm{div}}$ ablation (12-epoch robust-aug MNIST).}
\label{tab:ndiv}
\begin{tabular}{ccccccc}
\toprule
$n_{\mathrm{div}}$ & Params & Clean & Rot45 & Rot90 & Noise & Lat.\ (ms/512) \\
\midrule
2 & 342k & 0.988 & 0.967 & 0.834 & 0.935 & 14.5 \\
4 & 271k & 0.987 & 0.966 & 0.887 & 0.885 & 12.1 \\
8 & 253k & 0.977 & 0.965 & 0.930 & 0.838 & 15.5 \\
\bottomrule
\end{tabular}
\end{table}

\begin{table}[H]
\centering
\small
\caption{Clean-only training (no rotation aug): FasterNet vs p4 equivariant. Equivariant is slightly better at Rot90/TTA but both remain weak without augmentation---so aug+TTA stays essential.}
\label{tab:clean_eq}
\begin{tabular}{lcccc}
\toprule
Model (clean train) & Clean & Rot45 & Rot90 & Rot90+TTA \\
\midrule
FasterNet-clean & 0.993 & 0.588 & 0.127 & 0.659 \\
p4 Equivariant-clean & 0.989 & 0.436 & 0.174 & 0.698 \\
\bottomrule
\end{tabular}
\end{table}

\begin{figure}[H]
\centering
\safeimg[0.82\linewidth]{fig_mnist_vs_emnist_bars.png}
\caption{MNIST vs EMNIST Digits robustness.}
\end{figure}

\begin{figure}[H]
\centering
\safeimg[0.90\linewidth]{fig_mnist_multangle_0to9.png}
\caption{MNIST digits 0--9 under multiple rotation angles.}
\end{figure}

\begin{figure}[H]
\centering
\safeimg[0.88\linewidth]{fig_mnist_multicorr_0to9.png}
\caption{MNIST digits 0--9 under multiple corruptions (noise, blur, shift, occlusion, elastic)---not only rotation.}
\end{figure}

\begin{figure}[H]
\centering
\safeimg[0.95\linewidth]{fig_mnist_rot90_vs_tta_0to9.png}
\caption{Single-pass Rot90 vs rotation TTA on digits 0--9; TTA recovers many orientation mistakes.}
\end{figure}

\begin{figure}[H]
\centering
\safeimg[0.72\linewidth]{fasternet_aug_rotate_90_perclass_cm.png}
\caption{MNIST confusion matrix at $90^\circ$ rotation; errors concentrate on confusable pairs such as $6\leftrightarrow 9$.}
\end{figure}

\begin{figure}[H]
\centering
\safeimg[0.78\linewidth]{fig_ndiv_ablation.png}
\caption{$n_{\mathrm{div}}$ accuracy--latency trade-off.}
\end{figure}

\begin{figure}[H]
\centering
\safeimg[0.82\linewidth]{fasternet_aug_svhn_long_corruption_bars.png}
\caption{SVHN corruption bars after 24-epoch training.}
\end{figure}

\FloatBarrier

\section{What I Have Learnt from This AI Assignment}
\begin{itemize}
  \item \textbf{Prompting is the real coursework:} when I only said ``train a CNN,'' results looked like toy upright MNIST. When I insisted on PConv checks, multi-seed bars, EMNIST/SVHN and ``not only digit 3,'' the project became a pattern-recognition study.
  \item \textbf{Clean accuracy lies:} clean-only FasterNet is $\sim$99\% upright but $\sim$13\% at Rot90; robust aug + TTA recovers $\sim$80--94\%. That gap is the lesson, not the headline clean score.
  \item \textbf{Architecture vs augmentation:} training a p4 equivariant CNN \emph{without} rotation aug tests whether invariance can come from the model itself (Table~\ref{tab:clean_eq}). FasterNet still needs aug for strong Rot90, so I keep FasterNet+aug+TTA as the main story and equivariance as a calibrated baseline.
  \item \textbf{Mechanism + variance matter:} PConv identity/FLOPs, $n_{\mathrm{div}}$ ablation and mean$\pm$std stop an AI report from looking like a single lucky run.
  \item \textbf{Domains matter:} EMNIST writer style and SVHN house numbers punish ``MNIST-only'' narratives; longer SVHN training fixed a large under-training gap.
  \item \textbf{Process logs help:} \texttt{process\_log.md} is required and also made debugging AI iterations easier.
\end{itemize}

\section{Source Code Webpage Link}
Source code:\\[0.3em]
\url{https://github.com/AngusJai/CISC3024-AI-Assignment1}\\[0.3em]
{\small Includes \texttt{models/}, training/evaluation scripts, \texttt{scripts/}, \texttt{process\_log.md}, and \texttt{report.tex}. Datasets (\texttt{data/}) and \texttt{.venv/} are gitignored. A trained MNIST FasterNet-aug checkpoint is attached as a GitHub Release asset for demo convenience.}

\section{Conclusion}
FasterNet with Partial Convolution, directed end-to-end through Cursor, is a credible recent Deep CNN for robust digit pattern recognition: verified PConv behaviour, strong multi-condition results on MNIST/EMNIST/SVHN versus clean-only collapse, ablations (including clean equivariant vs clean FasterNet), and a public source link. The assignment taught me that supervising AI---constraints, verification and honest limits---is as important as the model itself.

\begin{thebibliography}{9}
\bibitem{chen2023fasternet}
J.~Chen et al.\ Run, Don't Walk: Chasing Higher FLOPS for Faster Neural Networks.\ CVPR, 2023.
\bibitem{vasu2023mobileone}
P.~K.~A.~Vasu et al.\ MobileOne: An Improved One Millisecond Mobile Backbone.\ CVPR, 2023.
\bibitem{hendrycks2019benchmarking}
D.~Hendrycks and T.~Dietterich.\ Benchmarking Neural Network Robustness to Common Corruptions and Perturbations.\ ICLR, 2019.
\bibitem{lecun1998gradient}
Y.~LeCun et al.\ Gradient-based learning applied to document recognition.\ Proc.\ IEEE, 1998.
\bibitem{cohen2017emnist}
G.~Cohen et al.\ EMNIST: an extension of MNIST to handwritten letters.\ 2017.
\bibitem{netzer2011reading}
Y.~Netzer et al.\ Reading Digits in Natural Images with Unsupervised Feature Learning.\ NIPS Workshop, 2011.
\bibitem{rath2023invariances}
M.~Rath and A.~Condurache.\ Deep Neural Networks with Efficient Guaranteed Invariances.\ AISTATS, 2023.
\end{thebibliography}

\end{document}
